\documentclass[11pt, letterpaper]{article}
\usepackage{amsmath, amssymb}
\usepackage{hyperref}
\usepackage{booktabs}
\usepackage{enumitem}
\usepackage{xcolor}
\usepackage{geometry}
\usepackage{fancyhdr}
\usepackage{titlesec}
\usepackage{tcolorbox}
\usepackage{graphicx}
\usepackage{parskip}
\usepackage{caption}
\geometry{margin=1in}

\definecolor{courseblue}{HTML}{002452}
\definecolor{coursegold}{HTML}{FABD0F}
\definecolor{coursered}{HTML}{B90E31}
\definecolor{lightblue}{HTML}{E8EEF7}
\definecolor{lightgold}{HTML}{FEF8E3}

\hypersetup{colorlinks=true,linkcolor=courseblue,urlcolor=coursered}
\tcbuselibrary{skins}
\newtcolorbox{guidancebox}{enhanced,colback=lightblue,colframe=courseblue,
  sharp corners,leftrule=4pt,left=6pt,right=6pt,top=4pt,bottom=4pt}
\newtcolorbox{notebox}{enhanced,colback=lightgold,colframe=coursegold,
  sharp corners,leftrule=4pt,left=6pt,right=6pt,top=4pt,bottom=4pt}

\pagestyle{fancy}\fancyhf{}
\fancyhead[L]{\small\textcolor{courseblue}{\textbf{CISC 473 --- Deep Learning | Fall 2026}}}
\fancyhead[R]{\small\textcolor{courseblue}{Midterm Report}}
\fancyfoot[C]{\small\thepage}
\renewcommand{\headrulewidth}{0.4pt}
\titleformat{\section}{\large\bfseries\color{courseblue}}{}{0em}{}[\titlerule]
\titleformat{\subsection}{\normalsize\bfseries\color{courseblue}}{}{0em}{}

\begin{document}

{\centering
  {\Large\textbf{CISC 473: Deep Learning --- Midterm Report}}\\[0.5em]
  {\large\textit{P20: Speech Emotion Recognition with Self-Supervised Pre-Trained Features}}\\[0.4em]
  \textbf{Group Members:}\\{}
  Nicholas Saloufakos \quad 20398428 \quad Project Lead\\{}
  Connor Zakary \quad 20419228 \quad Research Lead \\{}
  Kaan Sun \quad 20346401 \quad Engineering Lead \\{}
  \url{https://github.com/nsaloufakos/Speech-Emotion-Recognition}\\[0.3em]
  \today\\[0.4em]
  \rule{\linewidth}{0.4pt}
}

\section{Abstract}

Speech emotion recognition aims to identify the emotional state conveyed by speech. But differences between speakers can complicated generalization to unfamiliar voices. Using the the eight-class RAVDESS dataset, this project explores how acoustic features, pretrained speech representations, and pitch normalization affect emotion classification. We have implemented an experimental pipeline comprising audio pre-processing, an MFCC-based bidirectional LSTM, and frozen HuBERT and wav2vec 2.0 representations with linear classifiers. Evaluation uses leave-one-out speaker cross validation across 24 actors and three random seeds, with separate validation speakers for model selection. Preliminary results show that frozen HuBERT representations achieve a mean unweighted average recall of 59.9\% on original recordings. Following normalization of each recording’s estimated mean fundamental frequency to 200 Hz, HuBERT performance falls to 43.0\%, compared with 35.5\% for wav2vec 2.0 and 34.5\% for the MFCC–BiLSTM baseline under the same normalized condition. This decrease suggests that pitch normalization may remove useful emotional information or introduce distortions. Next, we plan to complete comparisons on original audio, examine representations across HuBERT layers, evaluate fine-tuning, and test pitch, temporal, and noise augmentation to assess robustness.

\section{Introduction and Problem Statement}

\begin{guidancebox}
\textbf{Guidance.} Brief recap of the problem and why it matters (1/2 page max).
Focus on what has changed since the proposal---have you refined the scope or identified
a clearer novel angle? Include the \textbf{research question} you are now testing.
\end{guidancebox}

[CONCISE PROBLEM DESCRIPTION AND RESEARCH QUESTION. NOTE ANY SCOPE CHANGES.]

\section{Related Work (Updated)}

\begin{guidancebox}
\textbf{Guidance.} An extended version of your proposal survey. Add papers you
discovered since the proposal. Include a comparative table or taxonomy where possible.
Make explicit connections to your own approach. Add at least 2 new papers beyond
the proposal. Target: $\approx$400--600 words.
\end{guidancebox}

[UPDATED LITERATURE REVIEW.]

\begin{center}
\begin{tabular}{llp{4cm}p{4cm}}
\toprule
\textbf{Method} & \textbf{Year} & \textbf{Key idea} & \textbf{Relation to our work} \\
\midrule
{[Method 1]} & [Year] & [Key idea] & [Relation] \\
{[Method 2]} & [Year] & [Key idea] & [Relation] \\
{[Method 3] }& [Year] & [Key idea] & [Relation] \\
\bottomrule
\end{tabular}
\end{center}

\section{Methodology}

\begin{guidancebox}
\textbf{Guidance.} Describe what you have implemented in detail. Include:
(1) architecture / model description with any modifications,
(2) training procedure, hyperparameters, and optimiser,
(3) data preprocessing pipeline,
(4) evaluation protocol (metric definitions and dataset splits).
Use equations where appropriate.
\end{guidancebox}

We implemented three baseline models for eight-class speech emotion recognition on RAVDESS: an MFCC-based bidirectional long short-term memory network (BiLSTM), a frozen HuBERT encoder with a linear classifier, and a frozen wav2vec~2.0 encoder with a linear classifier. Layer-wise probing was also performed on the HuBERT model to investigate how accessible emotion information is at different Transformer depths. All models use a shared training and evaluation pipeline.

\subsection{Model Architecture}

\textbf{MFCC + BiLSTM}

Our first approach extracts acoustic features from each audio file by calculating 40 Mel-Frequency Cepstral Coefficients (MFCCs) for each time frame. These MFCC features are then passed through a single-layer Bidirectional Long Short-Term Memory (BiLSTM) model with 64 hidden units in each direction. The final forward and backward hidden states are combined to create an 128-dimensional utterance representation which then is mapped to 8 logits using a fully connected layer. Both the BiLSTM and the classification layer are randomly initialised and trained from scratch.

\textbf{HuBERT + Linear Classifier}

Our second approach initalises a HuBERT model using the pretrained weights from the checkpoint \texttt{facebook/hubert-base-ls960} to extract features from each audio file. All pretrained backbone weights remain frozen and unchanged during training. The encoder is kept in evaluation mode and used only to extract acoustic representations. For each recording, the final HuBERT  transformer layer produces a sequence of 768-dimensional hidden states, then mean time pooling is used to create a fixed-length utterance representation. To this backbone we added a randomly initialized linear classification layer with 768 input features and eight outputs that maps the representation to a classification predictions. Only the weights and biases of this added layer are updated during training.

\textbf{Wav2vec 2.0 + Linear Classifier}

Our third approach is very similar to the second, but instead initializes a pretrained wav2vec 2.0 model from the \texttt{facebook/wav2vec2-base} checkpoint. This model remains frozen and is only used to extract features from the audio files. It's final transformer layer is mean-pooled over time and the resulting 768-dimensional vector is mapped to eight emotion logits using a linear layer. The pretrained backbone weights remain unchanged. Only the weights and biases of the added, randomly initialized 768-to-8 linear classification layer are trained.

\textbf{HuBERT Layer-wise Probing}

To perform layer-wise probing we use the frozen pretrained weights from \texttt{facebook/hubert-base-ls960}. We extract the outputs of its 12 Transformer layers. Each layer's hidden states are mean-pooled over time to produce a 768-dimensional utterance representation. A separate, randomly initialized 768-to-8 linear classifier is trained for each layer using identical dataset splits and training settings.

\subsection{Training Setup}
All three baseline models and layer-wise probing classifiers were trained using the same procedures and hyperparameters. The baseline models were trained twice, first with no data augmentation or normalization, and a second time with Mean-F0 Semitone Normalization.

\begin{tabular}{ll}
\toprule
\textbf{Hyperparameter} & \textbf{Value} \\
\midrule
Optimiser & Adam \\
Learning rate & $1 \times 10^{-3}$ \\
Batch size & 32 \\
Epochs / iterations & 30 per fold and seed \\
Loss function & Unweighted Cross-entropy \\
Data augmentation & None \\
Data Normalization & Mean-F0 Semitone Normalization \\
Validation Seed & 42 \\
Random seeds used & 42, 123, 999 \\
\bottomrule
\end{tabular}

The training pipeline uses leave-one-speaker-out (LOSO) evaluation across all 24 actors. In each fold, one actor is reserved for testing. Two of the remaining actors are selected for validation using a fixed split seed of 42, and the other 21 actors are used for training. This produces 1,260 training recordings, 120 validation recordings, and 60 test recordings per fold. Split assignments are saved and reused across models and training seeds.

\subsection{Dataset and Preprocessing}

The experiments use the audio-only speech subset of RAVDESS, with
eight emotion labels: neutral, calm, happy, sad, angry, fearful,
disgust, and surprised. The complete subset contains 1,440 recordings
from 24 actors. The class distribution is mildly imbalanced: neutral has 96 recordings, while each of the other seven emotions has 192. To account for this we report both UAR and WAR. A dataset index is constructed from the recording filenames to identify each recording's emotion, actor, intensity, statement, and repetition. 

Audio is loaded as mono and resampled to 16~kHz. Mean-F0 Semitone Normalization is then performed and it's output is placed in a separate folder, as to not overwrite the original dataset. For the MFCC baseline, 40 MFCCs are extracted per time frame using Librosa. The resulting sequences are padded within each batch and packed before entering the BiLSTM, so that padding does not affect its final hidden states. For HuBERT and wav2vec~2.0, waveforms are processed using the associated feature extractor before being passed through the frozen encoder. Each recording is encoded individually, and its hidden states are mean-pooled over time to obtain a fixed-length representation.

We compare original audio with a separately generated version processed using Mean-F0 Semitone Normalisation. The same actor splits are used in both conditions. Normalisation is applied to the waveform before extracting MFCCs or pretrained speech representations. For each recording, we estimate its mean fundamental frequency, $\bar{F}_0$, from voiced frames. We then calculate the constant pitch shift required to move this mean towards a shared target frequency, $F_{\mathrm{target}}$:
\[
    \Delta s
    =
    12\log_2\left(
        \frac{F_{\mathrm{target}}}{\bar{F}_0}
    \right),
\]
where $\Delta s$ is the pitch shift in semitones. The same shift is applied throughout the recording, raising recordings whose mean pitch is below the target and lowering those above it. The purpose is to reduce differences in average pitch between speakers while retaining relative pitch movements within each recording. The procedure is deterministic and operates on each utterance independently, without using emotion labels or statistics from other recordings by the same speaker. It can therefore be applied to a single recording from an unseen speaker.

\subsection{Evaluation Protocol}

The primary evaluation metrics we use are unweighted average recall (UAR) and weighted average recall (WAR). For $C=8$ classes, they are defined as:
\[
    \mathrm{UAR}
    =
    \frac{1}{C}\sum_{c=1}^{C}
    \frac{\mathrm{TP}_c}{\mathrm{TP}_c+\mathrm{FN}_c}
    \qquad
    \mathrm{WAR}
    =
    \sum_{c=1}^{C}\frac{N_c}{N}
    \frac{\mathrm{TP}_c}{\mathrm{TP}_c+\mathrm{FN}_c}.
\]
where $N_c$ is the number of test recordings in class $c$, $N$ is the total number of test recordings, $\mathrm{TP}_c$ is the number of true positives, and $\mathrm{FN}_c$ is the number of false negatives for emotion class $c$. Validation performance is measured after every epoch. Training runs for all 30 epochs, and then checkpoint with the highest validation unweighted average recall (UAR) is retained. This checkpoint is then reloaded for the test evaluation. Training loss, validation loss, and training and validation UAR are recorded for each epoch. For each training seed, predictions from all evaluated test folds are pooled before calculating UAR and WAR. We then report the
mean and sample standard deviation across the three seed-level scores, using seeds 42, 123, and 999.

\section{Preliminary Results}

\subsection{Baseline Comparison}

We completed experiments for all three baseline models on original audio and on audio processed using Mean-F0 Semitone Normalization with a target of 200~Hz. Each experiment used leave-one-speaker-out evaluation across 24 actors and three training seeds (42, 123, and 999). For each seed, test predictions from all folds were pooled before calculating UAR and WAR. Table~\ref{tab:midterm-baselines} reports the mean and sample standard deviation across the three seed-level scores.

\begin{center}
\small
\begin{tabular}{llcc}
\toprule
\textbf{Model} & \textbf{Audio condition} & \textbf{UAR (\%)} & \textbf{WAR (\%)} \\
\midrule
MFCC--BiLSTM & Original & $35.07 \pm 1.14$ & $37.31 \pm 1.25$ \\
MFCC--BiLSTM & F0-normalized & $34.51 \pm 2.16$ & $36.57 \pm 1.99$ \\
HuBERT + linear & Original & $59.92 \pm 0.95$ & $60.81 \pm 0.91$ \\
HuBERT + linear & F0-normalized & $42.99 \pm 0.71$ & $44.13 \pm 0.74$ \\
wav2vec 2.0 + linear & Original & $44.10 \pm 0.14$ & $45.23 \pm 0.04$ \\
wav2vec 2.0 + linear & F0-normalized & $35.47 \pm 0.52$ & $36.41 \pm 0.35$ \\
\bottomrule
\end{tabular}
\captionof{table}{Baseline performance on eight-class RAVDESS under 24-fold LOSO evaluation. Values are mean $\pm$ sample standard deviation over three seeds. Original-audio experiments contain 1,440 recordings; the normalized condition contains 1,432 because eight recordings were excluded during preprocessing.}
\label{tab:midterm-baselines}
\end{center}

HuBERT achieved the highest UAR and WAR among the three original-audio baselines. Normalization reduced mean UAR by 16.93 percentage points for HuBERT and 8.63 points for wav2vec~2.0. The MFCC--BiLSTM decrease was much smaller, at 0.56 points, relative to its variation across seeds. These results do not support an improvement from the current normalization procedure.

For published context, Pepino et al.\ (2021) reported $84.3 \pm 1.7\%$ average recall on RAVDESS using combined wav2vec~2.0 layers, a dense classifier, and per-speaker feature normalization.\footnote{L. Pepino, P. Riera, and L. Ferrer, ``Emotion Recognition from Speech Using Wav2vec 2.0 Embeddings,'' Interspeech, 2021. \url{https://arxiv.org/abs/2104.03502}.} Their experiment merged neutral and calm into seven classes and used a fixed actor split, whereas ours retains eight classes and uses LOSO. Their model and normalization also differ from ours. The published score therefore provides context rather than a matched reproduction target. A matched reproduction remains outstanding.

\subsection{HuBERT Layer-wise Probing}

Probing all 12 HuBERT Transformer layers on original audio showed that the final layer was not the strongest representation for our linear classifier. Selected results are shown in Table~\ref{tab:midterm-layers}.

\begin{center}
\begin{tabular}{lcc}
\toprule
\textbf{Transformer layer} & \textbf{UAR (\%)} & \textbf{WAR (\%)} \\
\midrule
1 & $48.16 \pm 0.98$ & $50.28 \pm 0.84$ \\
6 & $69.38 \pm 0.31$ & $70.14 \pm 0.24$ \\
7 & $68.45 \pm 0.48$ & $69.10 \pm 0.43$ \\
12 & $59.92 \pm 0.95$ & $60.81 \pm 0.91$ \\
\bottomrule
\end{tabular}
\captionof{table}{Selected frozen HuBERT probing results on original recordings, using the same LOSO protocol and three seeds. Layer 6 had the highest observed test UAR across the 12 layers.}
\label{tab:midterm-layers}
\end{center}

Layer 6 exceeded the final layer by 9.46 UAR points, suggesting that emotion information is more accessible to a linear classifier at intermediate depths in this setting. This is a descriptive probing result, rather than a validation-selected improvement: layer 7 had the highest mean validation UAR. Probing under normalization remains future work.

\subsection{Qualitative Analysis}

The saved predictions provide examples of the errors introduced in the normalized condition. Table~\ref{tab:midterm-examples} shows two recordings classified correctly by original-audio HuBERT but incorrectly by its normalized-audio counterpart. These examples illustrate individual failures and were selected for inspection; they do not represent all predictions.

\begin{center}
\small
\begin{tabular}{lccc}
\toprule
\textbf{Recording ID} & \textbf{True emotion} & \textbf{Original prediction} & \textbf{Normalized prediction} \\
\midrule
\texttt{03-01-03-01-01-01-01} & Happy & Happy & Surprised \\
\texttt{03-01-04-01-01-01-16} & Sad & Sad & Disgust \\
\bottomrule
\end{tabular}
\captionof{table}{Illustrative HuBERT test predictions for seed 42. Each prediction comes from the fold holding out the recording's actor.}
\label{tab:midterm-examples}
\end{center}

Across all three seeds, mean HuBERT recall for happy recordings fell from 44.97\% to 17.88\% after normalization. Sad was already difficult on original audio, with 29.34\% recall, compared with 82.47\% for calm; the most frequent original-audio error was sad being classified as calm. This indicates that aggregate performance hides substantial differences between emotions.

\subsection{Discussion}

The current results establish a working comparison pipeline and show that frozen pretrained representations, especially HuBERT's intermediate layers, can outperform our handcrafted-feature baseline. However, the proposed pitch normalization has not improved performance. Possible explanations include removing emotion-related differences in average pitch, unreliable F0 estimates, and artifacts introduced by pitch shifting. We have not isolated these causes, so we treat them as hypotheses. The unequal recording sets also limit the comparison. Further preprocessing validation and controlled experiments are needed before drawing broader conclusions about pitch normalization.

\section{Challenges and Adaptations}

The main conceptual challenge is separating speaker-related pitch differences from emotion-related pitch differences. Normalizing each utterance independently may remove both. Our initial hypothesis was that reducing average-pitch differences would help generalization to unseen speakers, but the present results do not support that hypothesis for the current implementation and target frequency.

A related technical challenge is validating the preprocessing. The implementation estimates F0 using pYIN, averages valid voiced frames, and applies a constant semitone shift. It does not save per-recording F0 estimates, confidence values, or applied shifts, and it does not check the resulting pitch after saving. Listening and sample re-estimation motivate further inspection, but re-estimation alone cannot establish a processing error: the estimator can classify different frames as voiced before and after transformation. We plan to inspect corresponding voiced frames, record confidence and shift diagnostics, and review audio quality before changing the method.

The normalization script also excludes recordings with fewer than three valid voiced F0 estimates. This reduced the dataset from 1,440 to 1,432 recordings, so the normalized experiments do not have identical recording counts to the original-audio experiments. We will document the exclusions and use matched recording sets in follow-up comparisons.

Relative to the proposal, baseline evaluation, original-audio layer probing, and an initial normalization comparison are complete. Fine-tuning scripts and augmentation code have been implemented, but evaluated fine-tuning and augmentation-ablation results are not yet available in the repository. We have therefore revised the remaining plan to prioritize preprocessing validation, followed by these outstanding comparisons. Any alternative pitch estimator, shift limit, or normalization statistic will be evaluated as a separate condition rather than presented as an already completed improvement.

\subsection{Updated Week 7--10 Plan}

\begin{center}
\begin{tabular}{lp{11.5cm}}
\toprule
\textbf{Week} & \textbf{Planned tasks} \\
\midrule
7 & Validate pitch estimates and processed audio; record preprocessing diagnostics and exclusions; prepare matched original/normalized recording sets. \\
8 & Evaluate fine-tuning and begin controlled augmentation comparisons, applying training augmentations only to training recordings. \\
9 & Complete prioritized comparisons and normalized HuBERT probing; analyze per-emotion errors and consolidate tables and figures. \\
10 & Finalize interpretation and limitations; complete the final report and presentation. \\
\bottomrule
\end{tabular}
\end{center}

The team will confirm task ownership and prioritize experiments according to available training time. Model settings will be chosen using validation data, with the same speaker-disjoint evaluation protocol retained for the final comparisons.

\section{Preliminary Novel Contribution (Optional Bonus)}

Our preliminary insight is that bringing different speakers into a common pitch range is less straightforward than the semitone-shift formula suggests. Matching an estimated mean F0 does not by itself standardize pitch range or make voices sound alike. To investigate this, we re-estimated F0 in original and saved normalized recordings using the same pYIN settings as the implementation. For Actor~3's happy recording \texttt{03-01-03-01-01-01-03}, the estimated mean changed from 175.25~Hz to 272.55~Hz, despite a normalization target of 200~Hz. However, the detector identified 36 voiced frames before processing and 56 afterward. Restricting both estimates to the 35 frames classified as voiced in both versions gave means of 175.97~Hz and 205.41~Hz, respectively. This illustrates how voiced-frame selection complicates validation; it does not by itself establish an incorrect pitch shift. Recognition performance also decreased after normalization. Together, these observations motivate further checks of pitch estimation, audio quality, and preservation of emotional expression.

\section{Contribution Declaration}

\begin{tabular}{lcc}
\toprule
\textbf{Member} & \textbf{Contributions (this period)} & \textbf{\%} \\
\midrule
{[Student 1]} & [e.g., Lit review, updated proposal section] & \% \\
{[Student 2]} & [e.g., Model implementation, training runs] & \% \\
{[Student 3] }& [e.g., Data pipeline, evaluation code] & \% \\
\bottomrule
\end{tabular}

\section{References}

\begin{enumerate}[label={[\arabic*]},leftmargin=*]
  \item [Author(s), Year. ``Title.'' \textit{Venue}. arXiv:XXXX.XXXXX]
  \item [\ldots]
\end{enumerate}

\end{document}
